% Eisenberg 2016: Visible Type Application - Key Examples
% Extracted from the extended version

\documentclass[11pt,a4paper]{article}

\usepackage{amsmath,amssymb,amsthm}
\usepackage{mathpartir}
\usepackage{geometry}
\usepackage{hyperref}
\usepackage{xcolor}

\geometry{margin=1in}

\title{Eisenberg 2016: Visible Type Application \\ Key Examples with Analysis}
\author{Extracted from ESOP 2016 (Extended version)}
\date{}

\begin{document}

\maketitle

\tableofcontents

\section{Introduction --- Motivating Examples}

\subsection{Example 1: Basic Visible Type Application (Line 67)}

\textbf{Code:}
\[
\text{id \,@\,Int}
\]

\textbf{System:} System HMV / System V

\textbf{Demonstrates:} Direct explicit type instantiation

\textbf{Why it works:}
\begin{itemize}
\item $\text{id}$ has type $\forall a.\, a \to a$
\item $@\text{Int}$ instantiates $a$ with $\text{Int}$
\item Result: $\text{Int} \to \text{Int}$
\end{itemize}

---

\subsection{Example 2: Proxy Workaround (Lines 153--176)}

\textbf{Without visible type application (requires Proxy):}
\[
\begin{array}{l}
\text{data Proxy } a = \text{Proxy} \\[6pt]
\text{normalizeProxy} :: \forall a.\, (\text{Show } a, \text{Read } a) \Rightarrow \text{Proxy } a \to \text{String} \to \text{String} \\
\text{normalizeProxy } x = \text{show } ((\text{read} :: \text{String} \to a)\, x) \\[6pt]
\text{normalizeExpr} = \text{normalizeProxy } (\text{Proxy} :: \text{Proxy Expr})
\end{array}
\]

\textbf{With visible type application (cleaner):}
\[
\begin{array}{l}
\text{normalize} :: \forall a.\, (\text{Show } a, \text{Read } a) \Rightarrow \text{String} \to \text{String} \\
\text{normalize } x = \text{show } (\text{read} \,@\,a\, x)
\end{array}
\]

\textbf{System:} Contrasts pre-2016 GHC vs. System HMV

\textbf{Demonstrates:} Eliminating proxy types

\textbf{Why it works:} Visible type application allows direct specification of type arguments without runtime proxy values.

---

\section{System HM vs System HMV}

\subsection{Example 3: Type with Multiple Variables --- Subsumption (Lines 526--532)}

\textbf{Code} (Subsumption examples $\leq_{\text{hmv}}$):

\textbf{These work:}
\[
\begin{array}{ll}
\forall \{a, b\}.\, a \to b \;\leq_{\text{hmv}}\; \forall \{a\}.\, a \to a & \text{-- Instantiates } b \text{ with } a \\[4pt]
\forall a, b.\, a \to b \;\leq_{\text{hmv}}\; \text{Int} \to \text{Int} & \text{-- Instantiates both} \\[4pt]
\forall a, b.\, a \to b \;\leq_{\text{hmv}}\; \forall a.\, a \to \text{Int} & \text{-- Instantiates } b \\[4pt]
\forall a, b.\, a \to b \;\leq_{\text{hmv}}\; \forall \{a, b\}.\, a \to b & \text{-- Generalizes both} \\[4pt]
\forall a, b.\, a \to b \;\leq_{\text{hmv}}\; \forall \{b\}.\, \text{Int} \to b & \text{-- Instantiates } a
\end{array}
\]

\textbf{These do NOT work:}
\[
\begin{array}{l}
\forall a, b.\, a \to b \;\not\leq_{\text{hmv}}\; \forall b.\, \text{Int} \to b \quad \text{-- Can't instantiate only } b \text{, need left-to-right} \\[4pt]
\forall \{a\}.\, a \to a \;\not\leq_{\text{hmv}}\; \forall a.\, a \to a \quad \text{-- Can't subsume generalized with specified}
\end{array}
\]

\textbf{System:} System HMV (Higher-order Hindley-Milner with visible types)

\textbf{Demonstrates:} Subsumption relation $\leq_{\text{hmv}}$

\textbf{Why some work:} HMV\_InstS rule allows instantiating a tail of specified variables left-to-right

\textbf{Why some fail:} Can't skip variables or mix generalized/specified incorrectly

---

\section{System V --- Syntax-Directed Algorithm}

\subsection{Example 4: Lazy Instantiation (Line 859)}

\textbf{Code:}
\[
(\text{let } x = (\lambda y.\, y :: \forall a.\, a \to a) \text{ in } x) \,@\,\text{Int} \;:\; \text{Int} \to \text{Int}
\]

\textbf{System:} System V

\textbf{Demonstrates:} Visible type application after let-binding

\textbf{Why it works:}
\begin{enumerate}
\item $\lambda y.\, y :: \forall a.\, a \to a$ has \textbf{specified} type $\forall a.\, a \to a$
\item $\text{let } x = \ldots \text{ in } x$ preserves the specified type
\item $x \,@\,\text{Int}$ can instantiate because type is specified (lazy instantiation)
\end{enumerate}

\textbf{Key insight:} If we eagerly instantiated at variable lookup, $x$ would be monomorphic and $@\text{Int}$ wouldn't work.

---

\section{System SB --- Higher-Rank with Bidirectional Checking}

\subsection{Example 5: Higher-Rank Function Arguments (Lines 918--924)}

\textbf{Without visible type application in lambda:}
\[
\begin{array}{l}
\text{let foo} :: (\forall a.\, a \to a) \to (\text{Int}, \text{Bool}) \\
\quad \text{foo} = \lambda f \to (f\, 3, f\, \text{True}) \\
\text{in foo id}
\end{array}
\]

\textbf{With visible type application in lambda:}
\[
\begin{array}{l}
\text{let foo} :: (\forall a.\, a \to a) \to (\text{Int} \to \text{Int}, \text{Bool}) \\
\quad \text{foo} = \lambda f \to (f \,@\,\text{Int}, f\, \text{True}) \\
\text{in foo id}
\end{array}
\]

\textbf{System:} System SB

\textbf{Demonstrates:} Higher-rank polymorphism + visible type application

\textbf{Why it works:}
\begin{itemize}
\item $f$ has type $\forall a.\, a \to a$ (higher-rank, polymorphic)
\item $f \,@\,\text{Int}$ explicitly instantiates before application
\item $f\, \text{True}$ implicitly instantiates via unification
\end{itemize}

\textbf{Rules involved:} SB\_DAbs (checking lambda), SB\_TApp (type application)

---

\subsection{Example 6: Mixed Instantiation Order (Lines 928--933)}

\textbf{Code:}
\[
\begin{array}{l}
\text{pair} :: \forall a.\, a \to \forall b.\, b \to (a, b) \\
\text{pair} = \lambda x\, y \to (x, y) \\[6pt]
\text{bar} = \text{pair } '\text{x}' \,@\,\text{Bool}
\end{array}
\]

\textbf{Result:} bar inferred to have type $\text{Bool} \to (\text{Char}, \text{Bool})$

\textbf{System:} System SB

\textbf{Demonstrates:} Partial instantiation, order matters

\textbf{Why it works:}
\begin{enumerate}
\item $\text{pair } '\text{x}'$ instantiates $a = \text{Char}$, now has type $\forall b.\, b \to (\text{Char}, b)$
\item $@\text{Bool}$ instantiates $b = \text{Bool}$
\item Final type: $\text{Bool} \to (\text{Char}, \text{Bool})$
\end{enumerate}

\textbf{Key point:} Type application can happen after some arguments are applied.

---

\subsection{Example 7: Two Subsumption Relations (Lines 1176--1190)}

\textbf{Weak subsumption ($\leq_b$):}
\[
\begin{array}{l}
\leq_b \forall a.\, a \to \text{Bool} \to \text{Bool} \quad\quad\quad \text{-- Can skip a top-level quantifier} \\[4pt]
\leq_b (\forall a.\, a \to a) \to \text{Bool} \quad\quad\quad \text{-- Contravariant instantiation} \\[4pt]
\text{Int} \to \forall a, b.\, a \to b \;\not\leq_b\; \text{Int} \to \forall b.\, \text{Bool} \to b \quad \text{-- Can't reorder in } \leq_b
\end{array}
\]

\textbf{Deep skolemization ($\leq_{\text{dsk}}$):}
\[
\begin{array}{l}
\text{Int} \to \forall a, b.\, a \to b \;\leq_{\text{dsk}}\; \text{Int} \to \forall b.\, \text{Bool} \to b \quad \text{-- } \leq_{\text{dsk}} \text{ CAN reorder} \\[4pt]
\text{Int} \to \forall a.\, a \to a \;\leq_{\text{dsk}}\; \forall a.\, \text{Int} \to a \to a \quad \text{-- Can instantiate spec vars in any order} \\[4pt]
(\text{Int} \to \forall b.\, b \to b) \to \text{Int} \;\leq_{\text{dsk}}\; (\forall a, b.\, a \to b \to b) \to \text{Int} \quad \text{-- Contravariant out-of-order inst.} \\[4pt]
\forall \{a\}.\, a \to a \;\leq_{\text{dsk}}\; \forall a.\, a \to a \quad \text{-- Generalized } \leq \text{ specified}
\end{array}
\]

\textbf{System:} System SB

\textbf{Demonstrates:} Two subsumption relations

\textbf{Why different:}
\begin{itemize}
\item $\leq_b$: Used at synthesis/checking boundary, must instantiate left-to-right
\item $\leq_{\text{dsk}}$: Used for checking against specified types, more flexible
\end{itemize}

---

\subsection{Example 8: Why Deep Skolemization is Necessary (Lines 1258--1308, Fig. 10)}

\textbf{Setup:}
\[
\text{Assume } \Gamma = x : \forall \{a\}.\, \text{Int} \to a \to a
\]

\textbf{Goal:} Type-check $(x :: \text{Int} \to \forall a.\, a \to a)$

\textbf{Without deep skolemization:}
\[
\begin{array}{l}
\Gamma \vdash_{\text{sb}} x \Rightarrow \text{Int} \to a \to a \quad \text{-- Instantiates } \forall \{a\} \\[4pt]
\text{Int} \to a \to a \;\not\leq\; \text{Int} \to \forall a.\, a \to a \quad \text{-- Can't subsume, } a \text{ is free!}
\end{array}
\]

\textbf{With deep skolemization:}
\[
\begin{array}{l}
\Gamma \vdash_{\text{sb}} x \Rightarrow \text{Int} \to a \to a \\[4pt]
\text{Int} \to a \to a \;\leq_{\text{dsk}}\; \text{Int} \to \forall a.\, a \to a \quad \text{-- Deep skolemization handles inner quantifiers}
\end{array}
\]

\textbf{System:} System SB

\textbf{Demonstrates:} Fig. 10 example

\textbf{Why it works:} SB\_DeepSkol rule (line 1219) pushes quantifiers inward before checking.

---

\section{Extended Examples (Appendix A)}

\subsection{Example 9: Type-Class Constraint with Proxy (Lines 1636--1744)}

\textbf{Before visible type application:}
\[
\text{assume} :: \forall a.\, \text{Proxy } c \to (c \Rightarrow a) \to a
\]

\textbf{Using it:}
\[
\text{the} :: \text{List } a \,\text{Unbranched} \to a \\
\text{the} = \text{assume } (\text{Proxy} :: \text{Proxy } (b \sim \text{Unbranched}))
\]

\textbf{After visible type application:}
\[
\begin{array}{l}
\text{assume} :: \forall a.\, (c \Rightarrow a) \to a \\[4pt]
\text{the} = \text{assume } @(b \sim \text{Unbranched})
\end{array}
\]

\textbf{System:} Real-world GHC extension

\textbf{Demonstrates:} Visible type application for type-class constraints

\textbf{Why it works:} $@(b \sim \text{Unbranched})$ explicitly specifies the equality constraint.

---

\subsection{Example 10: Dependently-Typed Programming (Lines 1744--1870)}

\textbf{Before (with 4 Proxy arguments):}
\[
\text{fact} :: \forall t\, f\, s\, b.\, \text{Sing } b \to \text{Proxy } t \to \text{Proxy } f \to \text{Proxy } s \to \text{Proxy } (\text{If } b\, t\, (f\, s)) \to \text{Proxy } (f\, (f\, s))
\]

\textbf{Using it:}
\[
\text{fact } (\text{sEval } se_0)\, (\text{Proxy} :: \text{Proxy } (\text{Eval } e_1))\, (\text{Proxy} :: \text{Proxy } (\text{Eval } e_2))\, (\text{Proxy} :: \text{Proxy } s)
\]

\textbf{After (with visible type application, only 1 Proxy):}
\[
\begin{array}{l}
\text{fact} :: \forall t\, f\, s\, b.\, \text{Sing } b \to \text{Proxy } (\text{If } b\, t\, (f\, s)) \to \text{Proxy } (f\, (f\, s)) \\[4pt]
\text{fact } @(\text{Eval } e_1)\, @(\text{Eval } e_2)\, @s\, (\text{sEval } se_0)
\end{array}
\]

\textbf{System:} Real-world GHC with GADTs

\textbf{Demonstrates:} Eliminating 3 out of 4 Proxy arguments

\textbf{Why it works:} Type application specifies $t$, $f$, $s$ directly.

---

\section{Summary Table: Examples by System}

\begin{center}
\begin{tabular}{|c|l|l|}
\hline
\textbf{Example} & \textbf{System} & \textbf{Key Feature Demonstrated} \\
\hline
1 & HMV/V & Basic visible type application \\
2 & HMV/V & Proxy elimination \\
3 & HMV & Subsumption relations \\
4 & V & Lazy instantiation \\
5 & SB & Higher-rank + visible types \\
6 & SB & Partial instantiation \\
7 & SB & Two subsumption relations \\
8 & SB & Deep skolemization necessity \\
9 & Real GHC & Type-class constraints \\
10 & Real GHC & Dependently-typed programming \\
\hline
\end{tabular}
\end{center}

---

\section{Key Takeaways}

\begin{enumerate}
\item \textbf{Lazy instantiation is crucial} --- without it, visible type application wouldn't work on let-bound variables

\item \textbf{System SB requires 5 mutually recursive judgments} --- our current implementation doesn't have this structure

\item \textbf{Two subsumption relations} ($\leq_b$ and $\leq_{\text{dsk}}$) handle different scenarios

\item \textbf{Deep skolemization} is necessary for higher-rank types with inner quantification

\item \textbf{Pattern matching} should preserve polymorphic types as ``specified'' for lazy instantiation
\end{enumerate}

\end{document}
